% =============================================================================
% METHODOLOGY - CITA Paper in ALKALI Format
% =============================================================================

\vspace{-2mm}
\section{Methodology}
\label{sec:methodology}

We present the CITA training methodology, organized into three phases: (1) setup and data preparation, (2) instruction-tuned formatting, and (3) contrastive training configuration.

\vspace{-2mm}
\subsection{Phase 0: Model and Infrastructure Setup}
\label{sec:phase0}

\begin{table}[h!]
\centering
\small
\begin{tabular}{@{}ll@{}}
\toprule
\textbf{Component} & \textbf{Specification} \\
\midrule
Base Model & Llama-3.1-8B~\cite{dubey2024llama} \\
Prompt Format & Chat-style (system/user/assistant) \\
Tokenizer & LLaMA-3 with BOS/EOS tokens \\
Hardware & NVIDIA A100-40GB \\
Training Data & PKU-SafeRLHF~\cite{ji2024pku} \\
\bottomrule
\end{tabular}
\caption{Training infrastructure and configuration.}
\label{tab:infrastructure}
\vspace{-4mm}
\end{table}

\textbf{Rationale.} We select Llama-3.1-8B for its strong instruction-following capabilities and accessibility. PKU-SafeRLHF provides high-quality preference pairs with explicit safety annotations, enabling evaluation of instruction-conditioned safety behaviors.

\vspace{-2mm}
\subsection{Phase 1: Dataset Formatting}
\label{sec:phase1}

We define two complementary data formats that enable instruction-conditioned training:

\paragraph{Instruction-Tuned Alignment (ITA) Format.}
Each training example explicitly includes an alignment instruction alongside the user query:

\begin{small}
\begin{verbatim}
<user>
### Alignment Instruction:
${alignment_instruction}
### User Prompt:
${user_prompt}
<assistant>
${expected_response}
\end{verbatim}
\end{small}

\paragraph{Contrastive ITA Format.}
For preference learning, we extend ITA to include rejected responses:

\begin{small}
\begin{verbatim}
### Rejected Response (Bad):
${rejected_response}
### Expected Response (Good):
${accepted_response}
\end{verbatim}
\end{small}

\paragraph{Example-Based Alignment (EBA) Format.}
Standard DPO format without explicit instructions for baseline comparison:
\begin{center}
\texttt{\{prompt, chosen, rejected\}}
\end{center}

\vspace{-2mm}
\subsection{Phase 2: Training Configuration}
\label{sec:phase2}

\begin{table}[h!]
\centering
\small
\begin{tabular}{@{}lcc@{}}
\toprule
\textbf{Parameter} & \textbf{SFT Stage} & \textbf{CITA Stage} \\
\midrule
Trainer & SFTTrainer & Custom DPO+KL \\
Epochs & 3--5 & 3 \\
Sequence Length & 2048 & 2048 \\
Learning Rate & 2e-5 & 5--7e-6 \\
LoRA Rank & 8 & 8 \\
$\beta$ (temperature) & --- & 0.1--0.12 \\
$\lambda_{\text{KL}}$ & --- & 0.0002--0.0005 \\
\bottomrule
\end{tabular}
\caption{Training hyperparameters for SFT and CITA stages.}
\label{tab:training_config}
\vspace{-4mm}
\end{table}

\textbf{Key Design Choice.} The CITA stage uses significantly lower learning rates (3--4$\times$ lower than typical DPO) because instruction-augmented sequences are 30--40\% longer, producing larger gradients. This finding emerged from hyperparameter optimization (Section~\ref{sec:hyperparams}).
